% Generated draft: requires mathematical and source review.
\documentclass[11pt]{article}
\usepackage[margin=25mm]{geometry}
\usepackage{fontspec}
\setmainfont{Latin Modern Roman}
\usepackage{amsmath,amssymb,mathtools,unicode-math}
\setmathfont{Latin Modern Math}
\usepackage{enumitem,needspace,xurl,hyperref}
\setlength{\parindent}{0pt}
\setlength{\parskip}{5pt}
\setlength{\emergencystretch}{3em}
\newcommand{\N}{\mathbb{N}}
\newcommand{\Z}{\mathbb{Z}}
\newcommand{\Q}{\mathbb{Q}}
\newcommand{\R}{\mathbb{R}}
\newcommand{\C}{\mathbb{C}}
\newcommand{\F}{\mathbb{F}}
\newcommand{\A}{\mathbb{A}}
\newcommand{\PP}{\mathbb{P}}
\newcommand{\E}{\mathbb{E}}
\newcommand{\eps}{\varepsilon}
\newcommand{\dd}{\,\mathrm{d}}
\begin{document}
\section*{UnsolvedMath ID: 30004693}
\textbf{Source problem code:} catalogue entry
\subsection{Definitions and mathematical statement}
\textbf{Scope and terminology.} The supplied record concerns mathematical classics across different cultural traditions, their availability to readers, their reading practices, and the implications of those practices for understanding mathematical argument and practice. It specifies no corpus of texts, cultural traditions, historical periods, reader populations, or criteria for assessing an explanatory answer. These terms therefore remain descriptive rather than formally defined mathematical objects.

\textbf{Recorded question.}
\begin{quote}
For mathematical classics in different cultural traditions, how were the texts made available to readers, how were they read, and what can these reading practices teach about the nature of mathematical argument and practice?
\end{quote}

\textbf{Unresolved formulation qualification.} This is a historiographical research question, not a specified mathematical proposition with defined domains, hypotheses, and quantifiers. The supplied registry marks its statement status as \texttt{unrecoverable} and supplies no clean statement. Its verification note reports that the workshop-level overview merges availability, reading-practice, and mathematical-argument themes retained separately in records \texttt{OWR-7155447-003} and \texttt{OWR-7155447-004}. The present record remains visible for provenance; the separate records' statements are not supplied here. Consequently, the question above is preserved as the recorded research prompt, and no unique formal mathematical assertion can be recovered from the supplied evidence.
\subsection{Short English statement}
The recorded question asks how readers in different cultural traditions obtained and read mathematical classics, and what those practices reveal about mathematical argument and practice. The registry identifies this entry as a merged workshop overview with an unrecoverable standalone statement. A precise mathematical proposition is therefore unavailable from the supplied record; this draft preserves the historical question and its limitations.
\subsection{Sources}
\begin{itemize}
\item Local UnsolvedMath registry, permanent ID \texttt{30004693}, source problem code \texttt{catalogue entry}, title \emph{Reading Practices and Arguments in Mathematical Classics}, in the \emph{Oberwolfach Reports Open Problems} collection. The supplied record lists proposed year 2021 and names no individual proposer.
\item Supplied report citation: \emph{Mathematics and its Ancient Classics Worldwide: Translations, Appropriations, Reconstructions, Roles} (2022), \emph{Oberwolfach Reports}, pp.\ 1347--1406, DOI \texttt{10.4171/owr/2021/26}. Supplied source URL and literature-source link, labeled ``Oberwolfach report'': \url{https://doi.org/10.4171/owr/2021/26}. The background describes Oberwolfach Reports as documentation of mathematical workshops and their open-problem discussions; no passage-level locator for this question is supplied.
\item Supplied statement-verification note: the official report is said to show that this overview merges themes retained separately in \texttt{OWR-7155447-003} and \texttt{OWR-7155447-004}; no separate clean statement is promoted, and the original record remains publicly visible for provenance. This is a report of the supplied registry assessment, not an independent verification of the report.
\item Supplied literature assessment, checked \textbf{2026-08-22}: status ``open''; the question is described as unresolved historiography concerning access to and reading of ancient mathematical texts, with workshop discussion rather than a settled explanatory result. This is supplied dated curation derived from the OWR material, not an independently verified assessment. The source expressly calls for independent verification of time-sensitive status and cited claims.
\end{itemize}
\end{document}
